% !TeX program = lualatex
\documentclass[11pt, a4paper]{article}

% --- EVRENSEL ÖNSÖZ BLOĞU ---
\usepackage[a4paper, top=2.5cm, bottom=2.5cm, left=2cm, right=2cm]{geometry}
\usepackage{fontspec}
\usepackage[turkish, bidi=basic, provide=*]{babel}

% Noto Sans font ayarları
\babelprovide[import, onchar=ids fonts]{turkish}
\babelprovide[import, onchar=ids fonts]{english}
\babelfont{rm}{Noto Sans}

% Liste simgeleri için
\usepackage{enumitem}
\setlist[itemize]{label=-}

% Matematik ve Tablo Paketleri
\usepackage{amsmath, amssymb, amsthm}
\usepackage{booktabs}
\usepackage{array}

% Bölüm numaralandırmasını kaldırmak için (isteğe bağlı)
\setcounter{secnumdepth}{0}

\begin{document}

\title{Depo Toplama Rotası Optimizasyonu: Matematiksel Model}
\author{}
\date{}
\maketitle

\section{Kümeler ve İndisler}

\begin{itemize}
    \item $i \in I$: Ürün artikel kodları kümesi, $I = \{1, \dots, 3388\}$
    \item $j_1 \in J_1$: Katlar, $J_1 = \{1, 2, 3, 4, 5, 6\}$
    \item $j_2 \in J_2$: Koridorlar, $J_2 = \{1, \dots, 27\}$
    \item $j_3 \in J_3$: Koridor tarafı (Sol: 0, Sağ: 1), $J_3 = \{0, 1\}$
    \item $j_4 \in J_4$: Sütunlar, $J_4 = \{1, \dots, 20\}$
    \item $j_5 \in J_5$: Raflar, $J_5 = \{1, 2, 3, 4, 5, 6, 7\}$
    \item $j_6 \in J_6$: THM (Taşıma Hazırlama Merkezi) kodu
    \item $j = (j_1, j_2, j_3, j_4, j_5, j_6)$: 6 boyutlu tam lokasyon indisi, $j \in J$
    \item $n, m = (j_1, j_2, j_4)$: 3 boyutlu fiziksel konum indisi, $n, m \in N$
\end{itemize}

\section{Parametreler}

\begin{itemize}
    \item $d_{n,m}$: $n$ konumu ile $m$ konumu arasındaki Manhattan mesafesi
    \item $s_{ij}$: $i$ ürününün $j$ lokasyonundaki mevcut stok miktarı
    \item $t_i$: $i$ ürününden toplanması gereken toplam miktar
    \item $w_1, w_2, w_3$: Amaç fonksiyonu ağırlıkları (Sırasıyla mesafe, kutu sayısı ve kat sayısı için)
\end{itemize}

\section{Karar Değişkenleri}

\begin{itemize}
    \item $x_{ij} \in \{0, 1\}$: $i$ ürününün $j$ lokasyonundan toplanıp toplanmama durumu
    \item $y_{ij} \geq 0$: $i$ ürününün $j$ lokasyonundan toplanan miktarı
    \item $u_n \in \{0, 1\}$: $n$ konumunun ziyaret edilip edilmeme durumu
    \item $v_{n,m} \in \{0, 1\}$: $n$ konumundan $m$ konumuna doğrudan gidilme durumu
    \item $p_n$: $n$ konumunun tur içindeki ziyaret sırası (MTZ kısıtı için)
    \item $z_{j_1} \in \{0, 1\}$: $j_1$ numaralı kata girilip girilmediği durumu
    \item $b_{j_6} \in \{0, 1\}$: $j_6$ numaralı THM kutusunun açılıp açılmadığı durumu
\end{itemize}

\section{Amaç Fonksiyonu}

\begin{equation}
    \min Z = w_1 \sum_{n \in N} \sum_{m \in N} d_{n,m} v_{n,m} + w_2 \sum_{j_6 \in J_6} b_{j_6} + w_3 \sum_{j_1 \in J_1} z_{j_1}
\end{equation}

\noindent \textbf{Terim Açıklamaları:}
\begin{itemize}
    \item \textbf{Terim 1:} Toplam yürünen Manhattan mesafesi.
    \item \textbf{Terim 2:} Toplam kullanılan (açılan) THM kutusu sayısı.
    \item \textbf{Terim 3:} Toplam girilen (aktifleşen) kat sayısı.
\end{itemize}

\section{Kısıtlar}

\begin{flalign}
    & y_{ij} \leq x_{ij} \cdot s_{ij} & \forall i \in I, \forall j \in J \label{cons:stok} \\
    & x_{ij} \leq b_{j_6} & \forall i \in I, \forall j \in J \label{cons:thm} \\
    & \sum_{j \in J} y_{ij} = t_i & \forall i \in I \label{cons:talep} \\
    & x_{ij} \leq u_n & \forall i \in I, \forall n \in N, \forall j \in J(n) \label{cons:ziyaret} \\
    & u_n \leq z_{j_1} & \forall j_1 \in J_1, \forall n \in N(j_1) \label{cons:kat} \\
    & \sum_{m \in N} v_{n,m} = u_n & \forall n \in N \label{cons:flow1} \\
    & \sum_{m \in N} v_{m,n} = u_n & \forall n \in N \label{cons:flow2} \\
    & p_n - p_m + |N| \cdot v_{n,m} \leq |N| - 1 & \forall n, m \in N, n \neq m \label{cons:mtz}
\end{flalign}

\noindent \textbf{Kısıt Açıklamaları:}
\begin{itemize}
    \item \textbf{(\ref{cons:stok})}: Bir lokasyondan toplanan ürün miktarı, o lokasyondaki stok miktarını geçemez.
    \item \textbf{(\ref{cons:thm})}: Eğer bir lokasyondan ürün toplandıysa, ilgili THM kutusu açılmış sayılır.
    \item \textbf{(\ref{cons:talep})}: Her ürün için toplanması gereken toplam miktar sağlanmalıdır.
    \item \textbf{(\ref{cons:ziyaret})}: Ürün toplama yapılan her lokasyonun bulunduğu fiziksel konuma uğranmış olmalıdır.
    \item \textbf{(\ref{cons:kat})}: Eğer bir konuma gidildiyse, o konumun bulunduğu kat ziyaret edilmiş sayılır.
    \item \textbf{(\ref{cons:flow1})-(\ref{cons:flow2})}: Akış korunumu kısıtları. Bir noktaya girildiyse oradan çıkılmalıdır.
    \item \textbf{(\ref{cons:mtz})}: Miller-Tucker-Zemlin (MTZ) alt tur eleme kısıtı. Rotaların tek bir kapalı döngü olmasını ve alt turların oluşmamasını sağlar.
\end{itemize}

\end{document}